\section{模型假设}\label{sec:assumptions}

为便于模型建立，在不改变题目基本运行规则的基础上，作如下假设：

\begin{enumerate}
    \item \textbf{忽略题目未给出的外部扰动及附加资源限制。}
    不考虑天气、风速、机械故障、临时空域变化等随机因素，同时不额外设置机场停机位、海上设施并发起降能力及加油设施服务能力等题目未规定的约束。

    \item \textbf{具有相同运输属性的人员具有等价性。}
    对起点、终点及其他相关属性完全一致的人员，可在求解阶段聚合为同一需求组，求解后再还原至具体人员；问题三还需保证组内人员的任务类型和时间要求一致。

    \item \textbf{问题三中的非临时任务均需完成。}
    将应急、增储上产和常规倒班任务视为必须满足的运输需求，仅临时任务允许根据剩余运输能力选择性安排。

    \item \textbf{问题三的时间要求通过时间可行架次列进行筛选。}
    候选方案仅保留满足人员最早上机、最晚到达及每日运行时段等时间规则的安排，不额外假设题目未规定的随机等待或延误；飞机返场后的30 min周转时间按题意计入具体飞机的连续占用约束。
\end{enumerate}
